% Penggerak mandiri untuk bab Deduksi Aksiomatik dalam pelokalan Bahasa
% Indonesia. Penggerak ini membatasi build pada OLP-0112--OLP-0125.

\documentclass[openany]{memoir}

% `cleveref` Open Logic belum memiliki opsi Indonesian. Babel tetap memuat
% Indonesian; `\ollanguage` bernilai english hanya selama cleveref dimuat.
\PassOptionsToPackage{english,indonesian}{babel}

\newcommand*{\olpath}{../..}
\newcommand*{\ollangid}{id}
\newcommand*{\ollanguage}{english}

\usepackage{mathpazo}
\usepackage{helvet}
\usepackage{courier}
\linespread{1.05}

\input{\olpath/sty/open-logic.sty}
\renewcommand*{\ollanguage}{indonesian}
\input{\olpath/sty/open-logic-defer.sty}

\includeenv{editorial}
\problemsperchapter
\let\cleardoublepage\clearpage
\pagestyle{plain}

% Pisahkan judul soal berturut-turut setelah daftar atau derivasi yang panjang.
\BeforeBeginEnvironment{probdeferred}{\par\medskip}
\BeforeBeginEnvironment{probdeferredtag}{\par\medskip}

\begin{document}

\selectlanguage{indonesian}

% OLP-0124 merujuk hasil sintaks dan semantik orde pertama yang belum
% termasuk dalam build bab terbatas ini. Entri ini mempertahankan identitas
% dan penomoran pembaca beku tanpa mengimpor prosa Inggris.
\makeatletter
\hypertarget{ol-id-forward-fol-syntax}{}
\protected@write\@auxout{}{\string\newlabel{fol:syn:sat:defn:satisfaction}{{16.{11}}{0}{Keterpenuhan}{ol-id-forward-fol-syntax}{}}}
\protected@write\@auxout{}{\string\newlabel{fol:syn:sat:defn:satisfaction@cref}{{[defn][11][16]16.{11}}{[1][0][]0}{}{}{}}}
\protected@write\@auxout{}{\string\newlabel{fol:syn:ass:prop:sentence-sat-true}{{16.{18}}{0}{}{ol-id-forward-fol-syntax}{}}}
\protected@write\@auxout{}{\string\newlabel{fol:syn:ass:prop:sentence-sat-true@cref}{{[prop][18][16]16.{18}}{[1][0][]0}{}{}{}}}
\protected@write\@auxout{}{\string\newlabel{fol:syn:ass:prop:sat-quant}{{16.{19}}{0}{}{ol-id-forward-fol-syntax}{}}}
\protected@write\@auxout{}{\string\newlabel{fol:syn:ass:prop:sat-quant@cref}{{[prop][19][16]16.{19}}{[1][0][]0}{}{}{}}}
\protected@write\@auxout{}{\string\newlabel{fol:syn:ext:prop:extensionality}{{16.{20}}{0}{Ekstensionalitas}{ol-id-forward-fol-syntax}{}}}
\protected@write\@auxout{}{\string\newlabel{fol:syn:ext:prop:extensionality@cref}{{[prop][20][16]16.{20}}{[1][0][]0}{}{}{}}}
\protected@write\@auxout{}{\string\newlabel{fol:syn:ext:cor:extensionality-sent}{{16.{21}}{0}{Ekstensionalitas untuk Kalimat}{ol-id-forward-fol-syntax}{}}}
\protected@write\@auxout{}{\string\newlabel{fol:syn:ext:cor:extensionality-sent@cref}{{[cor][21][16]16.{21}}{[1][0][]0}{}{}{}}}
\protected@write\@auxout{}{\string\newlabel{fol:syn:ext:prop:ext-formulas}{{16.{23}}{0}{}{ol-id-forward-fol-syntax}{}}}
\protected@write\@auxout{}{\string\newlabel{fol:syn:ext:prop:ext-formulas@cref}{{[prop][23][16]16.{23}}{[1][0][]0}{}{}{}}}
\protected@write\@auxout{}{\string\newlabel{fol:syn:sem:thm:sem-deduction}{{16.{30}}{0}{Teorema Deduksi Semantik}{ol-id-forward-fol-syntax}{}}}
\protected@write\@auxout{}{\string\newlabel{fol:syn:sem:thm:sem-deduction@cref}{{[thm][30][16]16.{30}}{[1][0][]0}{}{}{}}}
\makeatother

\olimport*[first-order-logic/axiomatic-deduction]{axiomatic-deduction}

\end{document}
